\clearpage
\fancyhead[L]{\small GRAHAMS PLANET.}
\SourceLeaf{482}
\begin{center}
\rule{0.55\linewidth}{0.4pt}\\[2pt]
Astronomische Nachrichten. Band XXVII. Nr. 639. S. 235. 236. 10. Juni 1848.\\[2pt]
\rule{0.55\linewidth}{0.4pt}
\end{center}
\begin{flushright}Göttingen 1848 Mai\raisebox{.35ex}{,} 22.\end{flushright}

Den neuen Planeten habe ich noch am Tag, an dem ich Ihre Nachricht erhielt, am Meridiankreis beobachtet, und zwei Tage später noch einmal. Beim ersten Mal war der Planet sehr lichtschwach. Beim zweiten Mal, als die Luft voller Moorrauch war, war er so lichtschwach, dass die Beobachtung nur wenig Vertrauen verdient. Trotzdem gebe ich hier beide an.

\begin{center}
\textsc{Graham}s Planet.\\[3pt]
\setlength{\tabcolsep}{4pt}
\begin{tabular}{l|r|r|r}
1848 & Göttinger m. Z. & Rektasz. & Dekl. \\\hline
Mai 9 & $11^{\rm u}\ 31^{\rm m}\ 3^{\rm s}95$ & $220^\circ\ 38'\ 6''6$ & $-11^\circ\ 57'\ 48''2$ \\
\hfill 11 & $11\ 21\ 14.7$ & $220\ 8\ 40.5$ & $-11\ 53\ 10.3$ \\
\end{tabular}
\end{center}

